\section*{6. \underline{消灭圈局部理论的应用}\\
\underline{(Picard--Lefschetz)}}

\subsection*{命题 6.1.}
\underline{设 $Z$ 为射影空间 $\mathbb P^r_k$ 中的超曲面, $r\geq2$, 其中 $k$ 为代数闭域. 设}
\[
i_o:D_o\lhook\joinrel\longrightarrow\mathbb P^r
\]
\underline{为一条直线的包含映射, 它仅在光滑部分与 $Z$ 相交, 且相交横截. 选取 $D_o-D_o\cap Z$ 的一个几何点 $\bar a$, 并以下式表示 $\pi_1(D_o-D_o\cap Z,\bar a)$ 的商群, 它分类 $D_o-D_o\cap Z$ 上在 $D_o\cap Z$ 的每个点处均驯分歧的平展覆盖:}
\[
\pi_1^{\mathrm{mod}}(D_o-D_o\cap Z,\bar a)
\]
\underline{则:}

\subsubsection*{6.1.1.}
\underline{典范同态}
\begin{equation}\tag{6.1.1.1}
i_{o*}:\pi_1^{\mathrm{mod}}(D_o-D_o\cap Z,\bar a)
\longrightarrow
\pi_1^{\mathrm{mod}}(\mathbb P^r-Z,\bar a)
\end{equation}
\underline{为满射.}

\subsubsection*{6.1.2.}
\underline{每选取一个几何点 $\bar x_o\in D_o\cap Z$ 以及 $D_o$ 中从 $\bar x_o$ 到 $\bar a$ 的一条“道路”, 都确定同态}
\begin{equation}\tag{6.1.2.1}
\prod_{\ell\ne\operatorname{char}(k)}\mathbb Z_\ell(1)
\longrightarrow
\pi_1^{\mathrm{mod}}(D_o-D_o\cap Z,\bar a)
\end{equation}
\underline{其像为 $\bar x_o$ 处的“那个”惯性群. 将此态射与 $i_{o*}$ 复合, 得到同态}
\begin{equation}\tag{6.1.2.2}
\prod_{\ell\ne\operatorname{char}(k)}\mathbb Z_\ell(1)
\longrightarrow
\pi_1^{\mathrm{mod}}(\mathbb P^r-Z,\bar a).
\end{equation}
\underline{若 $Z$ 不可约, 则所有这些同态彼此共轭.}

\underline{证明.}
设 $T$ 为 Grassmann 簇的一个子概形, 由 $\mathbb P^r$ 中满足以下条件的直线构成:
\begin{enumerate}[label=(\roman*)]
\item 通过 $\bar a$;
\item 仅在 $Z$ 的光滑部分与之相交, 且在那里相交横截.
\end{enumerate}
设 $\eta$ 为其泛点, 并设
\[
i_\eta:D_\eta\longrightarrow\mathbb P^r\otimes_k k(\eta)
\]
为相应的直线 (因而是“连接 $\bar a$ 与 $\mathbb P^r_k$ 的泛点”的直线). 再选取 $\eta$ 上的一个几何点 $\bar\eta$, 并以 $D_{\bar\eta}$ 表示相应的“几何”直线. 设 $o\in T$ 为与给定直线 $D_o$ 对应的点.
